% !PRINTED 493
\weberpage{493}
\weq{13}{\begin{aligned}
n &= 11,\quad & c &= \infty,\ 0,\ 1,\ 2,\ 3,\ 4,\ 5,\ 6,\ 7,\ \ 8,\ 9,\ 10,\\
&& c' &= \infty,\ 2,\ \ 3,\ 4,\ 5,\ 6,\ 7,\ 8,\ 9,\ 10,\ 0,\ \ 1,\\
&& c'' &= 0,\ \infty,\ 10,\ 5,\ 7,\ 8,\ 2,\ 9,\ 3,\ \ 4,\ 6,\ \ 1,\\
&& c''' &= 1,\ 10,\ \ 0,\ 4,\ 6,\ 5,\ 8,\ 7,\ 9,\ \ \ 2,\ 3,\ \infty.
\end{aligned}}

As roots of the resolvents of degree 7 and of degree 11 we may, following \S\,82, take in each case one of the two following functions:

\weq{14}{\begin{aligned}
n &= 7,\\
w_\nu &= \frac{(v_\infty - v_\nu)(v_{\nu+1}-v_{\nu+3})(v_{\nu+2}-v_{\nu+6})(v_{\nu+4}-v_{\nu+5})}{i\sqrt{7}\,u^4},\\
w'_\nu &= \frac{(v_\infty - v_\nu)(v_{\nu+1}-v_{\nu+5})(v_{\nu+2}-v_{\nu+3})(v_{\nu+4}-v_{\nu+6})}{-i\sqrt{7}\,u^4},
\end{aligned}}

\weq{15}{\begin{aligned}
n &= 11,\quad w_\nu =\\
&\frac{(v_\infty-v_\nu)(v_{\nu+1}-v_{\nu+2})(v_{\nu+3}-v_{\nu+6})(v_{\nu+4}-v_{\nu+8})(v_{\nu+5}-v_{\nu+10})(v_{\nu+9}-v_{1+7})}{i\sqrt{11}\,u^6},\\
w'_\nu &=\\
&\frac{(v_\infty-v_\nu)(v_{\nu+1}-v_{\nu+6})(v_{\nu+3}-v_{\nu+7})(v_{\nu+4}-v_{\nu+2})(v_{\nu+5}-v_{\nu+8})(v_{\nu+9}-v_{\nu+10})}{-i\sqrt{11}\,u^6},
\end{aligned}}
% [sic] German source prints the last factor as $(v_{\nu+9}-v_{1+7})$; the index pattern would require $v_{\nu+7}$.

in which $\sqrt{7}$ and $\sqrt{11}$ are to be taken positive.

There now emerges, first of all, the following remarkable difference between the two cases.

By the two permutations $(c')$, $(c'')$ the $w_\nu$ are, in the case $n = 7$, only permuted among themselves, namely in the following manner:

\[
\begin{aligned}
& w_0,\ w_1,\ w_2,\ w_3,\ w_4,\ w_5,\ w_6,\\
(c') \quad & w_2,\ w_3,\ w_4,\ w_5,\ w_6,\ w_0,\ w_1,\\
(c'') \quad & w_0,\ w_2,\ w_1,\ w_6,\ w_4,\ w_5,\ w_3,
\end{aligned}
\]

while [according to (13)] in the case $n = 11$ the quantities $w$ at the same time change their sign under the permutation $(c')$, so that the following permutations occur:

\[
\begin{aligned}
& w_0,\ w_1,\ w_2,\ w_3,\ w_4,\ w_5,\ w_6,\ w_7,\ w_8,\ w_9,\ w_{10},\\
(c') \quad & -w_2,\ -w_3,\ -w_4,\ -w_5,\ -w_6,\ -w_7,\ -w_8,\ -w_9,\ -w_{10},\ -w_0,\ -w_1,\\
(c'') \quad & w_0,\ w_9,\ w_6,\ w_3,\ w_5,\ w_4,\ w_2,\ w_8,\ w_7,\ w_1,\ w_{10},
\end{aligned}
\]

and it is the same with $w'$.

From this it now follows (\S\,54) that the $w_\nu$ are roots of an equation of degree 7 or of degree 11:

\weq{16}{w^n + A_1 w^{n-1} + A_2 w^{n-2} + \cdots + A_n = 0}
%CONT>>
